% ---------------------------------------------------------------------
%  Chapitre 5 — BLOC 3 (esquissé, mais résultats réels disponibles)
% ---------------------------------------------------------------------
\chapter{Bloc 3 --- relations inter-actifs}
\label{chap:bloc3}

Ce bloc étudie les séries de rendements sectoriels~: degré de couplage,
relations d'avance/retard, et prévisibilité du marché par son passé.

\section{Corrélations sectorielles}
% Résultat réel : corrélation moyenne hors-diagonale ~ 0,70
La corrélation moyenne hors-diagonale entre secteurs vaut environ
\textbf{0,70}~: le marché évolue largement « d'un bloc ». Les secteurs les
plus couplés sont \emph{Industrials} et \emph{Financials}~; les plus
diversifiants, \emph{Energy} et \emph{Utilities}.

\section{Causalité de Granger}
\label{sec:b3-granger}
% [À DÉVELOPPER] Définition : le passé de A améliore-t-il la prévision
% de B ? Correction de Bonferroni sur 110 paires.
Sur 110 paires ordonnées (retards 1--3 jours, correction de Bonferroni),
\textbf{48} sont significatives~\cite{granger1969}, \emph{Financials}
apparaissant comme meneur principal.

\begin{limites}
La très grande taille d'échantillon (\(\approx 4\,132\) jours) confère au
test une puissance telle que « significatif » n'implique pas
« exploitable » (gains de prévision faibles, retards courts). De plus,
« Granger-cause » signifie \emph{prédictif}, non causal au sens réel.
\end{limites}

\section{ACP des rendements --- le facteur marché}
% Résultat réel : PC1 = 73 % ; PC2 (8 %) = axe risk-on/risk-off
La première composante capte \textbf{73~\%} de la variance, tous les
secteurs de même signe~: c'est le \emph{facteur marché} (risque
systématique). La deuxième composante (8~\%) oppose secteurs défensifs et
cycliques.

\section{ARIMA sur le rendement de marché}
\label{sec:b3-arima}
% [À DÉVELOPPER] Test ADF (stationnarité), ACF/PACF, ARIMA(1,0,1).
Le prix est non stationnaire (marche aléatoire, ADF \(p\approx1\)), le
rendement est stationnaire (\(p\approx10^{-26}\))~; les autocorrélations
sont quasi nulles. Un \(\mathrm{ARIMA}(1,0,1)\)~\cite{boxjenkins1970} ne
dégage pas de structure exploitable.

\begin{conditions}
Travailler sur les \textbf{rendements} (stationnaires), jamais sur les
prix bruts~; vérifier la stationnarité (ADF) avant toute modélisation
ARIMA. % point à défendre à l'oral
\end{conditions}
